% ==================== PREÁMBULO ====================
% Archivo para centralizar la configuración de paquetes y estilos del documento.

% ----- IDIOMA Y CODIFICACIÓN -----
\usepackage[T1]{fontenc}
\usepackage[utf8]{inputenc}
\usepackage[spanish,english]{babel}

% ----- GEOMETRÍA Y TIPOGRAFÍA -----
\usepackage[letterpaper,margin=2.2cm]{geometry}
\usepackage{mathptmx} % Fuente tipo Times para formato académico/IEEE
\usepackage{setspace}
\setstretch{1.05}
\usepackage{parskip}
\usepackage{enumitem}
\setlist{nosep}

% ----- PAQUETES MATEMÁTICOS -----
\usepackage{amsmath,amssymb,amsfonts}

% ----- FIGURAS Y GRÁFICOS -----
\usepackage{graphicx}
\usepackage{xcolor}
\usepackage{float}
\usepackage{tikz}
\usetikzlibrary{mindmap, positioning, arrows.meta, shapes.geometric}

% ----- TABLAS -----
\usepackage{booktabs}
\usepackage{tabularx}
\usepackage{array}
\usepackage{multirow}
\usepackage{multicol}
\usepackage{threeparttable}
\usepackage{longtable}

% ----- BIBLIOGRAFÍA Y REFERENCIAS -----
\usepackage[numbers]{natbib}
\usepackage{xurl}      % Permite cortar enlaces largos en varias líneas
\usepackage{url}
\usepackage{hyperref}
\hypersetup{
  colorlinks=true,
  linkcolor=black,
  citecolor=black,
  urlcolor=blue,
  breaklinks=true
}
\Urlmuskip=0mu plus 2mu

% ----- CÓDIGO Y COMANDOS -----
\usepackage{listings}
\lstdefinestyle{terminal}{
  basicstyle=\ttfamily\small,
  frame=single,
  breaklines=true,
  columns=fullflexible,
  backgroundcolor=\color{gray!8},
  rulecolor=\color{gray!45},
  xleftmargin=0.5em,
  xrightmargin=0.5em
}
\lstset{style=terminal}

% ----- AJUSTES TIPOGRÁFICOS -----
\usepackage{etoolbox}
\hbadness=10000
\vbadness=10000
\tolerance=10000

% ----- COMANDOS PROPIOS -----
\newcommand{\evidencia}[2]{%
\begin{figure}[H]
\centering
\fbox{\begin{minipage}[c][0.18\textheight][c]{0.92\textwidth}
\centering
\textbf{#1}\\[0.4em]
\small Inserte aquí el recorte de pantalla: \texttt{#2}
\end{minipage}}
\caption{#1}
\end{figure}}
